\chapter{Technologies}
\label{chap:Technologies}
In this chapter we describe the technologies we have used for this project in detail. 
An overview of these technologies is given here:

\begin{itemize}
	\item Qt, a platform-independent framework for display and audio output.
	\item \gls{MIDI} to read input from instruments like keyboards.
	\item The libraries RtMIDI and libUSB to read MIDI from USB devices.
	\item We read scores from GP5 files.
	\item For audio synthesis we used SoundTouch.
	\item KissFFT provides simple frequency analysis functions for the audio analysis.
\end{itemize}

Qt is a cross-platform application framework written in \cpp from \emph{Digia}. It provides functions for GUI design via QML (see \ref{sec:QML}) and a signal-/slot-concept detailed next.

\section[Qt Signal-/Slot-concept \textsuperscript{Kugler}]{Qt Signal-/Slot-concept}
\label{sec:Signal-/Slot}

The Qt framework provides a signal/slot-concept to allow communication between classes in a thread-safe manner. Classes that inherit from {\tt QObject} can define signals and slots.

\subsection{Signals}

Signals are events other objects can listen to. 
They are defined in the class header after the {\tt signal} keyword.

\begin{lstlisting}
	class SignalExample : public QObject
	{
		Q_OBJECT
		
		public:
			// some variables
			
		signal:
			void mySignal(int someNumber);
	}
\end{lstlisting}

This class has a signal {\tt mySignal} with an integer argument. It can be emitted in a function with the emit keyword like this:

\begin{lstlisting}
	emit mySignal(5);
\end{lstlisting}

\subsection{Slots}

Similar to signals we can define slots using the {\tt slot} keyword. Slots can be specified public or private by using the corresponding keyword.

\begin{lstlisting}
	class SlotExample : public QObject
	{
		Q_OBJECT
		
		public:
			// some variables
			
		public slot:
			void mySlot(int someNumber);
	}
\end{lstlisting}

Slots have a function body that can be triggered by signals. They are defined like normal function implementations.

\begin{lstlisting}
	void SlotExample::mySlot(int someNumber) 
	{
		// do something with someNumber
		cout << someNumber << endl;
	}
\end{lstlisting}

\subsection{Connections}

To use the signals and slots we just created we can connect them by using the {\tt connect} function. 

\begin{lstlisting}
	SignalExample *signalObject = new SignalExample();
	SlotExample *slotObject = new SlotExample();

	connect(
		signalObject, &SignalExample::mySignal, 
		slotObject, &SlotExample::mySlot);
\end{lstlisting}

Now {\tt mySlot} will be called every time {\tt mySignal} is emitted. A signal can be connected to multiple slots and a slot can be connected to multiple signals. 

This mechanism provides a way for objects to communicate with each other and to write highly reusable code. 

\section[QML \textsuperscript{Novotny}]{QML}
\label{sec:QML}
QML is a markup language for defining user interfaces. It is part of the QtQuick framework. QML is loosely based on JSON. The basic structure of the language is declarative. JavaScript can be used inside a QML file which allows including imperative code in the user interface definition. This approach allows implementing both appearance and logic of the user interface in QML rather than \cpp. This section gives an overview of the basic functionalities of QML. Emphasis is given to features relevant to the application.
\subsection{Basics}
\paragraph{QML Components}
QML Components are specified by their type followed by a pair of braces. A QML component may contain child objects. These child components are defined within the braces. Furthermore every component has a set of properties that can be read and modified by other components. Values can be assigned to the properties between the two braces.
\paragraph{Properties}
The values of properties can be either primitive types or QML components. Basically there are two ways to set the value of property.
\begin{itemize}
\item{
	\textbf{Property Bindings}
	Whenever a property is assigned a JavaScript expression inside the definition of the component a property binding is created. This means that the value of the property will be updated automatically if a property included in the expression is updated. This allows specifying property values in a declarative way.
	\begin{lstlisting}
	  Text {
	      text: model.modelData.getTitle()
	      color: standardFontColor
	      font.bold: highlighted
	      font.pixelSize: standardFontPixelSize * 3
	      elide: Text.ElideRight
	      width: scoreDataItem.width  - 2 * marginInItem
	      font.family: headerFont
	  }
	\end{lstlisting}
	A number of property bindings are created in the code examples above. For example:
	\begin{itemize}
		\item {The {\tt text}-property is bound to the result of the function {\tt getTitle} of the property {\tt model.modelData}. The function {\tt getTitle()} will be evaluated when the property {\tt modelData} is updated.}
		\item{The {\tt color}-property of the {Text-} component is bound to the property {\tt standardFontColor} of a parent component.}
	\end{itemize}
}
\item {
	\textbf{Property Assignment}
	The value of a property can be set manually using JavaScript code. Manually setting the value of a property removes any existing bindings of the property.
	\begin{lstlisting}
		Rectangle {
		  Component.onCompleted: {
		      width = parent.width
		  }
		}
	\end{lstlisting}
	In the code example above the {\tt width} -property of the rectangle is set to the current value of the {\tt width} -property of its parent. The property will not be updated automatically if the parent component is updated.
}
\end{itemize}
\subsection{User Input}
\subsubsection{Handling Mouse Events}
QML provides an invisible component called {\tt MouseArea} which inherits {\tt Item}. {\tt MouseAreas} are positioned within other components in order to implement mouse handling. {\tt MouseArea} provides a number of signals such as:
\begin {itemize}
	\item{{\tt onEntered}}
	\item{{\tt onExited}}
	\item{{\tt onClicked}}
	\item{{\tt onDoubleClicked}}
\end{itemize}
\begin{lstlisting}
	Rectangle {
	    width: 100
	    height: 100

	    MouseArea {
	        anchors.fill: parent
	        onClicked: {
	            // Do Something
	        }
	    }
	}
\end{lstlisting}
The position of the mouse cursor can be retrieved via the properties {\tt mouseX} and {\tt mouseY}.
{\tt MouseArea} allows modifying the appearance of the mouse cursor inside the component.
\subsubsection{Handling Key Events}
All QML components provide a property {\tt Keys} to support key handling. This property provides the signals {\tt onPressed} and {\tt onReleased}.
	\begin{lstlisting}
		Rectangle {
	     	Keys.onPressed: {
		         if (event.key === Qt.Key_Return) {
		             event.accepted = true;
		             // Do something
		         }
	     	}
		}
	\end{lstlisting}
Note that by default the event will be forwarded to the parent element. This can be avoided by \textit{consuming} the event which is achieved by setting its {\tt accepted} -property to {\tt true}.
\subsection {Dynamically loading Components}
\paragraph{Loader}
The {\tt Loader} element allows dynamic loading of a subtree of QML objects. {\tt Loader} inherits {\tt Item} and can therefore be used in the same way. \\
This is used to load other QML objects on demand rather than when the application is started.  This leads to serious performance improvements when multiple complex objects need to be loaded at different points in time. 
\paragraph{}
The {\tt Loader} component can load a QML file. This has been used for implementing the application´s menu structure. \\
Alternatively it can load an existing QML object. This has been used in the implementation of  {\tt SelectableListView}, a powerful, flexible {\tt ListView} implementation and will be explained at a later point.
If the source of the Loader is changed the previously loaded Item is destroyed and resources are freed.
\subsection {ListView}
{\tt ListView}s in QML display data provided by data models.

% http://qt-project.org/doc/qt-5/qtquick-modelviewsdata-modelview.html
% \includegraphics[width=0.5\textwidth]{model_view}

They are highly flexible in terms of what data is displayed and how the data is displayed.
\subsubsection{Data Models}
What data is displayed is determined by data models.
These models can be defined in three ways.
\begin {itemize}
	\item {
		\textbf{QML, by inheriting ListModel} \\
		A ListModel in QML is a source of data. It consists of {\tt ListElements} which can be added dynamically to the list model. The attributes of the ListElements are defined dynamically when creating the records. \\
		These attributes are referred to as data roles. The following code example shows the definition of the list model behind the Display Option Control in the Game Start menu. The contents of the ListModel are specified explicitly inside its definition. The ListElements consist of a single role called \textit{name}.
		\begin{lstlisting}
			ListModel {
               id: displayOptionsModel
               ListElement { name: "Piano Roll"  }
               ListElement { name: "Score" }
               ListElement { name: "Tabulature"}
          	}
		\end{lstlisting}
	}
	\item {
		\textbf{\cpp, as a collection of string values} \\
		This has been used for representing lists of scalar data such as the instrumental tracks contained in a score.
	}
	
	\item {
		\textbf{\cpp, as a collection of QObjects}
		This has been used for representing lists of complex data such as scores and sections inside scores.	
	}
\end {itemize}
\subsubsection {Delegates}
Delegates represent the elements of a view visually. Each record in a view´s model is visualised by the template defined in the delegate. The data of the record can be accessed and modified by its delegate. The delegate may bind to the data roles of its associated record. If the name of a data role coincides with the name of a property of the delegate it can be accessed using the qualified model name instead. If the record is scalar it´s data may be accessed by using the property modelData. 
This feature helps making delegates reusable. 

\subsection{Communication between \cpp and QML}
Due to the tight integration of the Qt Meta Object System in QML passing data to QML objects takes little effort in most situations. Properties, methods and signals of a {\tt QObject} can be accessed from QML.

\subsubsection{Context Properties}
Context properties allow passing data obtained by \cpp code to QML components. Context properties affect all QML objects that are displayed by an {\tt ApplicationView} object. 

The following code examples show how a {\tt QObject} is passed to an {\tt ApplicationView} object as a context property and then used by a QML object.
\begin{lstlisting}
	class DataClass : public QObject
	{
	    Q_OBJECT
	
	    Q_PROPERTY(int number READ getNumber WRITE setNumber NOTIFY numberChanged)
	}
\end{lstlisting}
\begin{lstlisting}
	QQuickView view;
	view.show();
	QPointer<DataClass> data = new DataClass();
	view.rootContext()->setContextProperty("data", data);
\end{lstlisting}

\begin{lstlisting}
	Text {
		anchors.fill: parent
	    text: data.number
	}
\end{lstlisting}

\subsubsection{Making Data Accessible}
Objects passed to QML need to be instances of a class inheriting {\tt QObject}.
Basically there are two ways to make the data of an object available to QML objects:

\begin{itemize}
	\item{Properties}
	\item{Functions}
\end{itemize}

\paragraph{Properties}
In order to be accessible for QML objects class attributes can be wrapped in {\tt Propertie}s.
A property is a wrapper for class attributes. It is associated with a getter method,  an optional setter method and an optional signal that is emitted when the value of the property is changed.
Properties are defined using the {\tt Q\_PROPERTY} macro.

The following code example shows an extract of the definition of a class that uses properties to expose data to QML objects.
\begin{lstlisting}
	class Note : public QObject
	{
	    Q_OBJECT
	
	    Q_PROPERTY(int time READ getTime WRITE setTime NOTIFY timeChanged)
	    Q_PROPERTY(int duration READ getDuration WRITE setDuration NOTIFY durationChanged)
	    Q_PROPERTY(int pitch READ getPitch WRITE setPitch NOTIFY pitchChanged)
	
	private:
	    int time; // number of miliseconds since the song started, 0 if unknown
	    int duration; // 0 if the note is still playing
	    int pitch; // as specified by the MIDI standard, -1 if unknown, 0 is a valid pitch in MIDI

	public:
	    int getTime() const;
	    int getDuration() const;
	    int getPitch() const;
	
	public slots:
	    void setTime(int value);
	    void setDuration(int value);
	    void setPitch(int value);
	
	signals:
	    void timeChanged(int value);
	    void durationChanged(int value);
	    void pitchChanged(int value);
	};
\end{lstlisting}

\paragraph{Functions}
When used in front of the definition of a function the macro {\tt Q\_INVOKABLE} registers a function with the meta-object system. This allows invoking the function from a QML context. 
\subsubsection{Data Type Conversion}
When data is passed from a \cpp context to a QML context it needs to be converted to JavaScript data types. The QML engine will automatically convert most \cpp data types. For example a {\tt QString} object will be converted to a JavaScript string object when passed to a QML object and vice-versa.

\paragraph{Conversion of Custom Types}
Custom Types can be converted to JavaScript types if:
\begin {itemize}
	\item {the class extends {\tt QObject} and}
	\item {the custom type is registered with the QML type system.}
\end {itemize}


\section[Display methods in Qt \textsuperscript{Novotny}]{Display methods in Qt}
Qt provides multiple different ways to display graphics on the screen. Each has different advantages.

\subsection{Inheriting QQuickItem}
Custom components can be created using \cpp by extending the class {\tt QQuickItem}.
{\tt QQuickItem}s are painted using the scene graph approach (explained below).
This approach to drawing custom components has only been introduced with Qt 5. The official documentation lacks examples and detailed explanations. Refer to the presentation \textit{Integrating QML with C++} \cite{qtquickpainting1} by \textit{Digia Plc.} for a brief introduction to the topic.
\subsubsection{Qt Quick Scene Graph}
Traditionally, drawing of custom components in Qt has been based on imperative painting. Imperative painting causes a large amount of draw calls which has a negative affect on the drawing performance.

The scene graph concept has been introduced with Qt 5. The scene graph consists of nodes that are defined before the drawing process is started. This allows a number of optimizations and reduces the number of draw calls significantly.

\subsubsection{Nodes}
Nodes are parts of the scene graph. They represent objects that will be drawn. Node can be containers for other nodes. 

All node classes extend {\tt QSGNode}.
Subclasses used in the project will be discussed in detail.

\paragraph{Memory Managment}
The flag {\tt QSGNode::OwnedByParent} passes ownership of a node to its parent node. The nodes will be deleted as soon as the scene graph is destroyed. Setting this flag is highly recommended in most situations.

Similarly, material and geometry (refer to the explanation below) should be owned by the corresponding node in most situations. This is achieved by using the flags {\tt QSGNode::ownsMaterial} and {\tt QSGNode::ownsNode}.

\paragraph{QSGGeometryNode}
The {\tt QSGGeometryNode} class is the most important node class for drawing custom components. Many other classes such as {\tt QSGSimpleRectNode} extend {\tt QSGGeometryNode}.

A {\tt QSGGeometryNode} primarily consists of 
\begin {itemize}
	\item {Geometry}
	\item {Material}
\end {itemize}

The geometry of a node defines its position, shape and size. The simplest way of specifying the geometry of a node is defining its vertices.

The material defines how a node is filled.

The following code example shows how a line-shaped {\tt QSGGeometryNode} is created.
This example demonstrates how ownership flags are used to ease memory managment.

\begin{lstlisting}
QSGNode* QQI_Based_Visuals::getLineNode(qreal x0, qreal y0, qreal x1, qreal y1, QColor color) {
    QSGGeometryNode* lineNode = new QSGGeometryNode();
    QSGGeometry* lineGeometry = new QSGGeometry(QSGGeometry::defaultAttributes_Point2D(), 2);
    lineGeometry->setLineWidth(2);
    lineGeometry->setDrawingMode(GL_LINE_STRIP);
    lineNode->setGeometry(lineGeometry);
    QSGFlatColorMaterial* material = new QSGFlatColorMaterial();
    material->setColor(color);
    lineNode->setMaterial(material);

    QSGGeometry::Point2D *vertices = lineGeometry->vertexDataAsPoint2D();
    vertices[0].x = x0;
    vertices[0].y = y0;
    vertices[1].x = x1;
    vertices[1].y = y1;

    lineNode->setFlag(QSGNode::OwnedByParent);
    lineNode->setFlag(QSGNode::OwnsGeometry);
    lineNode->setFlag(QSGNode::OwnsMaterial);
    return lineNode;
}

\end{lstlisting}

\paragraph{QSGSimpleRectNode}
{\tt QSGSimpleRectNode} inherits {\tt QSGGeometryNode}. It allows creating rectangle-shaped nodes while simplifying the definition of the geometry. 
{\tt QSGSimpleRectNode} offers a constructor accepting a {\tt QRectF} object.

\begin{lstlisting}
	QRectF noteRectangle = QRectF(left, top, width, height);
	QSGSimpleRectNode* newRectNode = new QSGSimpleRectNode(noteRectangle, Qt::black);
\end{lstlisting}

\paragraph{QSGSimpleTextureNode}
{\tt QSGSimpleTextureNode} simplifies drawing textured objects. A {\tt QSGSimpleTextureNode} primarily consists of 
\begin{itemize}
	\item{texture}
	\item{bounding rectangle}
\end{itemize}

\begin{lstlisting}
	QSGSimpleTextureNode *digitNode = new QSGSimpleTextureNode();
	digitNode->setTexture(digitTextures[digits[i]]);
	digitNode->setRect(left + i * size, top, size, size);
	node->appendChildNode(digitNode);
\end{lstlisting}

\subsubsection{Painting Process}
{\tt QQuickItem} provides the function {\tt update} which causes repainting of the item. 
When the {\tt update} function is called the virtual function {\tt updatePaintNode} is called prior to the rendering process. 

The function {\tt updatePaintNode} receives a {\tt QSGNode} containing the contents of the scene graph at the time the function is called. Th
is allows a more efficient painting process in case only minor changes are made to the graph.

\subsection{Qt Quick Canvas Painting}
The QML {\tt Canvas} component allows painting custom components using Java Script code. 
{\tt Canvas} provides an empty area to draw on.

The {\tt Canvas} component implements the \textit{W3C HTML 5 Canvas} specification. This allows porting HTML 5 canvas components to QML with little effort. 

The article \textit{Qt Quick Painting using Canvas Item} \cite{qtquickpainting2} provided by \textit{Nokia} contains examples and best practices for using the canvas-based drawing approach.

\subsubsection{HTML5 Canvas Specification}
The canvas specification defines a set of drawing commands.
The most important commands include:
\begin{itemize}
	\item{{\tt beginPath} creates a new {\tt path}}
	\item{{\tt moveTo}}
	\item{{\tt line} adds a line-shaped path defined by the supplied coordinates}
	\item{{\tt rect} adds a rectangle-shaped path defined by the supplied coordinates}
	\item{{\tt text} adds a path consisting of text in a specific font}
	\item{{\tt stroke} draws the path}
\end{itemize}

\subsubsection{Canvas-Based Drawing in QML}
The {\tt Canvas} element provides the signal {\tt onPaint} which is emitted whenever repainting of the component is requested. 

The following example is an extract of the JavaScript-based implementation of the {\tt TabatureView} component. It demonstrates the canvas-based drawing process. Note that the component has been reimplemented using the {\tt QQuickItem}-based drawing approach due to performance issues.

\begin{lstlisting}
Canvas {
	id: canv
    anchors.fill: parent
    onPaint: {
        doTheDrawing(canv.context);
    }
    
    function drawLine(ctx, xFrom, yFrom, xTo, yTo, color, strength) {

        ctx.moveTo(xFrom, yFrom);
        ctx.lineTo(xTo, yTo);
        ctx.strokeStyle = color;
        ctx.lineWidth = strength;
        ctx.stroke();
    }
    
    function doTheDrawing(ctx) {
    
        for(var i = 1; i <= tabulatureView.numOfStrings(); i++) {
            drawLine(ctx, 0,  getStringLevel(i), canv.width, getStringLevel(i), "#000000", 1);
        }
 	}
}
\end{lstlisting}

As opposed to the {\tt QQuickItem}-based drawing approach the {\tt Canvas} painting approach does not make use of a scene graph which negatively affects performance when painting complex components. Performance problems occur especially in real-time applications.

\section[The MIDI protocol \textsuperscript{Kugler}]{The MIDI protocol}
\gls{MIDI} is a protocol to describe time-coded events. It is commonly used to transfer music notes from electronic instruments to computers and synthesizers. An important part of the project is to let people play with real instruments. With \gls{MIDI} we can read what keys the user plays from MIDI instruments like keyboards.

\subsection{Overview}
\gls{MIDI} serially transmits messages with a fixed size of 8 bit at a fixed rate. The first bit of a message determines whether the next 7 bits are a status code or data. The first 3 bit of the status code define the type of the message, the next 4 bit the channel. Therefore \gls{MIDI} supports up to 16 different channels. Devices can be set up to only listen to a specific channel. A status message is followed by up to two data messages, defining parameters of the status.

The most important messages are 
\begin{description}
	\item[Note-On] \hfill \\ Activates a note and contains information about pitch and volume.
	\item[Note-Off] \hfill \\ Deactivates the note.
	\item[Control Change] \hfill \\ Changes one of the 128 user-defined parameters. Most synthesizers allow configuration of how these values are interpreted.
	\item[Program Change] \hfill \\ Sets the program number for this channel. These define which instrument/synthesizer should play the notes.
\end{description}

MIDI can be transmitted over MIDI cables or encapsulated in other data streams over USB, FireWire, Ethernet, etc.
Beat the Score only supports MIDI over USB as standard MIDI connectors are rare nowadays. For the communication with MIDI devices libUSB and RtMidi have been used (see section \ref{sec:libusb and RtMidi}).

\subsection{General MIDI specification}
Since MIDI itself does not give any meaning to the program and control numbers the general midi standard has been created. The general MIDI specification defines for each of the 128 program numbers an instrument. Those range from different pianos, guitars and strings to synthetic sounds. For percussion the channel 10 is reserved and a percussion instrument is assigned to each pitch number. 


\section[GP5 Format \textsuperscript{Kugler}]{GP5 Format}
\label{sec:GP5}
To support many play techniques we decided to use the Guitar Pro 5 format, a proprietary file format developed by \emph{Arobas Music}. It is specialized for guitars but supports other instruments as well. 

This section uses the vocabulary described in chapter \ref{chap:Basic Musical Terms}.

Other than MIDI, it stores multiple additional informations, including but not limited to
\begin{itemize}
	\item number of strings on the guitar and tuning,
	\item how a string is picked,
	\item whether the note has been played as a harmonic note and if, what kind of harmonic note,
	\item what fingers are used to grip the an accord.
\end{itemize}

This information is used to display and detect notes more accurately for guitars.

Another big advantage is that songs in the GP5 format can be found in vast amounts on the internet.

The file consists of five parts:
\begin{enumerate}
	\item Meta information
	\item Song data
	\item Bars
	\item Track
	\item Beats
\end{enumerate}

\subsection{Meta information}

This contains information about the song and the score:
\begin{itemize}
	\item song title and subtitle
	\item name of the interpret
	\item name of the album
	\item name of the composer
	\item name of the transcriber
	\item page setup
\end{itemize}

\subsection{General song information}

Information about the actual song: 
\begin{itemize}
	\item tempo of the song
	\item signature
	\item the General MIDI instrument number for each track
\end{itemize}
		
\subsection{Bars}

This part contains an entry for every bar of the song and stores information about
\begin{itemize}
	\item signature changes
	\item information about repetition
	\item names and positions of the sections
\end{itemize}

\subsection{Track}

For every instrument in the song, additional information is stored. If the instrument is a guitar this sections contains how many strings are used and how they are tuned.

\subsection{Beat}

The actual notes for every bar and track. Informations stored here are:
\begin{itemize}
	\item length of the beat
	\item pitch (or pitches if an accord is played) of the beat
	\item information about bending
	\item linking to the next note
\end{itemize}

\section[SoundTouch \textsuperscript{Neumayer}]{SoundTouch}
\label{sec:SoundTouch}
To pitch missing audio samples from existing Wave files ({\tt .wav}) Beat The Score uses the SoundTouch library. In this section we describe how to use the SoundTouch library for pitching audio samples.
\subsection{Usage}
\begin{enumerate}
\item Create a {\tt SoundTouch} object and set predefined settings. In this example we choose to use the much faster quick seeking algorithm in SoundTouch's internal tempo changer routine and disable it's anti-aliasing filter. Changing the default settings results in dramatically faster pitching on ARM processors as CPU utilization is lowered.
\item Read a sound file (RIFF WAVE format) into memory and store a pointer to its data as {\tt tmpBuf} as well as the size of the {\tt char[]} behind that pointer. Note, the {\tt f.seek(44);} instructs the {\tt QFile} object to skip the WAVE file's RIFF header, it's format section as well as the data's header signature and the data length field, which results in a size of 44 bytes, according to the Wikipedia article about the RIFF WAVE format (http://de.wikipedia.org/wiki/RIFF\_WAVE).
\item Set by how many semi-tones the audio sample should be pitched. A positive value indicates pitching the sample up where as a negative value indicates pitching the sample down.
\item Put the to-be-pitched audio sample into SoundTouch's internal buffer and flush remaining samples from the buffer if available.
\item Create a new block of memory for storing the result of the pitching process. In the example the pointer {\tt newSample} points to this memory block. Since we have configured the SoundTouch library (at compile time) to use integer samples {\tt SAMPLETYPE} is a typedef to a \cpp {\tt short} type which has a size of 16 bits. In contrast, a \cpp {\tt char} has a size of 8 bits. This is the reason for both casting between {\tt SAMPLETYPE}s and {\tt char}s as well as multiplying SoundTouch's internal sample size by 2 for the new audio sample memory. The size of these types depends on the processor architecture but it applies to the main target architectures, x86 and ARM. Afterwards we pull the pitched data from SoundTouch's internal memory buffer into our newly created {\tt newSample}.
\end{enumerate}

\begin{lstlisting}
// Step 1
SoundTouch st;
st.setChannels(1);
st.setSampleRate(44100);
st.setSetting(SETTING_USE_QUICKSEEK, 1);
st.setSetting(SETTING_USE_AA_FILTER, 0);

// Step 2
QFile f(QString("test.wav"));
f.open(QIODevice::ReadOnly);
f.seek(44);
QByteArray barr = f.readAll();
f.close();
char * tmpBuf = barr.data();
int tmpBufSize = barr.size();

// Step 3
st.setPitchSemiTones(1);

// Step 4
st.putSamples((SAMPLETYPE*)tmpBuf, tmpBufSize);
st.flush();

// Step 5
int numSamples = st.numSamples();
int newSampleSize = numSamples * 2; // numSamples = number of 2 byte samples
char *newSample = new char[newSampleSize];
st.receiveSamples((SAMPLETYPE*)newSample, numSamples);
st.flush();
\end{lstlisting}


\section[KissFFT \textsuperscript{Kugler}]{KissFFT}
\label{sec:KissFFT}
The KissFFT library provides an implementation of the \gls{FFT} for frequency analysis. Here we describe how to use the library. See chapter \ref{chap:Audio analysis} for how we use the library to detect notes in audio signals.

\subsection{Usage}
To transform a signal from time domain into frequency domain we follow these steps:
\begin{enumerate}
	\item Create a {\tt kiss\_fft\_cfg} with the size of the signal. This will calculate all parameters for the given size. Most of the time you will always analyse the same number of samples so this object can be reused. Note that the size of the signal has to be a power of two.
	\item Write the signal into a {\tt kiss\_fft\_cpx} array. Hence we are dealing with real numbers and not complex ones, we just set the complex part of the entries to 0.
	\item Call the {\tt kiss\_fft} function with the configuration from step 1, the {\tt kiss\_fft\_cpx} array and a second {\tt kiss\_fft\_cpx} array for the result. Note that in-place calculation is not supported. If the same array is provided as input and output the function will create another buffer and write it back when the calculation finished.
\end{enumerate}

\begin{lstlisting}
	int fftSize = 1024;
	float *signal = new float[fftSize];

	// Step 1
	kiss_fft_cfg fftConfig = kiss_fft_alloc(fftSize, 0, NULL, NULL);
	
	// Step 2
	kiss_fft_cpx *complexSignal = new kiss_fft_cpx[fftSize];
	for (int i=0; i<fftSize; i++)
	{
		complexSignal[i] = (kiss_fft_cpx){signal[i], 0};
	}
	
	// Step 3
	kiss_fft_cpx *frequencySpectrum = new kiss_fft_cpx[fftSize];
	kiss_fft(fftConfig, complexSignal, frequencySpectrum);
	
	// frequencySpectrum contains now the result
	
	// cleanup
	delete complexSignal;
	delete signal;
	delete frequencySpectrum;
\end{lstlisting}

\subsection{Interpretation of the result}
The result contains the amplitude and phase for each sine component of the signal. We can calculate the index at which a given frequency can be found with
$$ i = f \frac{N}{samplerate} $$
where $i$ is the index in the result array, $f$ the frequency we are looking for, $samplerate$ the samplerate of the signal and $N$ the number of samples in the signal.

Because of the Nyquist-Shannon sampling theorem the frequencies above $\frac{N}{2}$ contain the same values as the lower half mirrored. They are not interesting for us.

The length of the complex entries describe the amplitude of the sine waves and the orientation describe the phase. Thus we can calculate the values using the trigonometry functions provide by {\tt math.h}.
\begin{lstlisting}
	float amplitude = hypotf(frequencySpectrum[index].r, frequencySpectrum[index].i);
	float phase = atan2f(frequencySpectrum[index].r, frequencySpectrum[index].i);
\end{lstlisting}


\section[libusb and RtMidi \textsuperscript{Neumayer}]{libusb and RtMidi}
\label{sec:libusb and RtMidi}

\subsection{RtMidi}
On desktop systems the way to handle MIDI signals happens through the RtMidi library. 
RtMidi uses the operating systems audio subsystem to get a list of connected MIDI devices and receive signals from them. On Windows RtMidi uses the Microsoft Multimedia MIDI API while on Linux it uses the Advanced Linux Sound Architecture API. The use of APIs for the respective platform is set at compile time of the RtMidi library.

\subsection{libusb}
As one of our goals was to make Beat The Score work on multiple platforms using the same code base we had to make sure that giving input to the application worked on all supported platforms in a similar way. To accomplish that on Android, we had to switch the library for receiving MIDI data as most Android devices do not have ALSA MIDI support enabled at the kernel level.
\\
The Android version of Beat The Score uses libusb instead of RtMidi to read data from the connected USB MIDI input device. We use a specially modified version of libusb which is forked from github.com/SpecLad/libusb-android. The forked version is available at github.com/beatthescore/libusb. The Android version of libusb allows the use of file descriptors to open the USB device through the Linux kernel. Since Android uses a fairly standard Linux kernel it works the same as other typical Linux systems where devices are represented as virtual files on the filesystem. To open the file descriptor of the USB device we use Java code in the form of broadcast receivers which asks the Android framework to open the file and return the file descriptor. The Java code is mandatory as there are no documented APIs that allow gaining the required permissions using purely native code.
\\
The way Android allows applications to access USB devices is by reading in the application's manifest file ({\tt AndroidManifest.xml}) and checking if the application requests the permission to access USB devices. On Android the user permits applications the required permissions at install time. The following code snippet specifies the request for the permission.
\begin{lstlisting}
<uses-feature android:name="android.hardware.usb.host"/>
\end{lstlisting}

When the application is started it has to register a so-called broadcast receiver that gets triggered when certain events happen on the system. An {\tt IntentFilter} is required to make the application only react to events that are relevant.
\begin{lstlisting}
IntentFilter filter = new IntentFilter(JNIUtils.ACTION_USB_PERMISSION);
registerReceiver(JNIUtils.UsbReceiver, filter);
\end{lstlisting}

This registers a receiver that only should be triggered when a USB device is connected.
\\
When the broadcast receiver is triggered it checks whether the user actually allowed the application to open the device. The following broadcast receiver resides in the {\tt JNIUtils} Java class.
\begin{lstlisting}
public static final BroadcastReceiver UsbReceiver = new BroadcastReceiver() {

public void onReceive(Context context, Intent intent) {
  if (ACTION_USB_PERMISSION.equals(intent.getAction())) {
    synchronized (context) {
      UsbDevice device = (UsbDevice)intent.getParcelableExtra(UsbManager.EXTRA_DEVICE);
	  if (intent.getBooleanExtra(UsbManager.EXTRA_PERMISSION_GRANTED, false)) {
        if(device != null) {
          UsbDeviceConnection conn = usbManager.openDevice(device);
          int descriptor = conn.getFileDescriptor();
          setConnected(device.getVendorId(), device.getProductId(), endPointAdr, descriptor);
        }
      } else {
        Log.d("BTS", "permission denied for device " + device);
      }
    }
  }
}
};
\end{lstlisting}

The call to {\tt setConnected()} actually is a JNI call to native code.

\begin{lstlisting}
extern "C" {
  JNIEXPORT void JNICALL Java_at_beatthescore_JNIUtils_setConnected(JNIEnv * env, jobject obj,
    jint vendorId, jint productId, jint endpoint, jint descriptor);
  JNIEXPORT void JNICALL Java_at_beatthescore_JNIUtils_setConnected(JNIEnv * env, jobject obj,
    jint vendorId, jint productId, jint endpoint, jint descriptor)
  {
    usbVendorid = (int)vendorId;
    usbProductid = (int)productId;
    usbDescriptor = (int)descriptor;
    usbEndPoint = (int)endpoint;
  }
}

\end{lstlisting}
This exposes the native \cpp function to the Java-part of the Android application through JNI (Java Native Interface). 
The following values have to be provided to {\tt setConnected()} for the native code to allow using the USB device.

\begin{description}
\item[Vendor and product identifier:] \hfill \\ These values are used to get the right {\tt libusb\_device} object.
\item[The USB device's end point:] \hfill \\ As USB devices can have multiple end points for data transmission the application needs to open the end point that provides the host with MIDI data.
\item[The file descriptor:] \hfill \\ The file descriptor for the USB device (provided by the operating system kernel and passed on to the application by the Android framework)
\end{description}

In order to catch this broadcast event, the application has to register broadcast receivers for dealing with attaching and detaching USB devices to the Android host.

\begin{lstlisting}
public static final BroadcastReceiver ConnectReceiver = new BroadcastReceiver() {

  public void onReceive(Context context, Intent intent) {
    if(intent.getAction().equals(UsbManager.ACTION_USB_DEVICE_ATTACHED))
    {
      UsbDevice device  = (UsbDevice)intent.getParcelableExtra(UsbManager.EXTRA_DEVICE);
      if(device != null) {
        for(int i = 0; i < device.getInterfaceCount(); i++) {
          UsbInterface tmpInterface = device.getInterface(i);
          for(int j = 0; j < tmpInterface.getEndpointCount(); j++) {
            UsbEndpoint endPoint = tmpInterface.getEndpoint(j);

            if(endPoint.getDirection() == 128) { //Direction from device to host
              endPointAdr = endPoint.getAddress();
              PendingIntent mPermissionIntent;
              mPermissionIntent = PendingIntent.getBroadcast(context, 0, new Intent(ACTION_USB_PERMISSION), 0);
              usbManager.requestPermission(device, mPermissionIntent);
            }
          }
        }
      }
    } else if(intent.getAction().equals(UsbManager.ACTION_USB_DEVICE_DETACHED)) {
      setConnected(-1, -1, -1, -1);
    }
  }
};
\end{lstlisting}

As with the {\tt JNIUtils.UsbReceiver} the {\tt JNIUtils.ConnectReceiver} has to be registered to react to broadcast events.

\begin{lstlisting}
connectFilter.addAction(UsbManager.ACTION_USB_DEVICE_ATTACHED);
connectFilter.addAction(UsbManager.ACTION_USB_DEVICE_DETACHED);
registerReceiver(JNIUtils.ConnectReceiver, connectFilter);
\end{lstlisting}

Opening the USB device is done by getting the right {\tt libusb\_device} object in a simple loop followed by opening the device using the provided file descriptor:
\begin{lstlisting}
libusb_device_handle *dev_handle;
libusb_context *ctx = NULL;
libusb_device **devs;
ssize_t cnt;

libusb_init(&ctx);
cnt = libusb_get_device_list(ctx, &devs);
libusb_device *selectedDevice = NULL;
for(i = 0; i < cnt; i++) {
  libusb_device_descriptor desc;
  int r = libusb_get_device_descriptor(devs[i], &desc);
  if (desc.idVendor == this->androidGlue->getVendorID() &&
      desc.idProduct == this->androidGlue->getProductID()) {
    selectedDevice = devs[i];
  }
}
libusb_open2(selectedDevice, &dev_handle, this->androidGlue->getDescriptor());
\end{lstlisting}

Reading data from the devices end point is done by calling {\tt libusb\_bulk\_transfer()}.
\begin{lstlisting}
int ret = libusb_bulk_transfer(dev_handle, this->androidGlue->getEndpoint(), data, 4, &nBytes, 3000);
\end{lstlisting}
The return value of the call allows the application to check whether the data transfer worked as intended. 
{\tt Note} objects will not be generated from the data when the transfer function returns one of the two following errors.
\begin{itemize}
\item LIBUSB\_ERROR\_TIMEOUT
\item LIBUSB\_ERROR\_OVERFLOW
\end{itemize}

\section[OpenSSL \textsuperscript{Neumayer}]{OpenSSL}
\label{sec:OpenSSL}
The license files that allow the use of Beat The Score are generated on the backend of the server side using the OpenSSL cryptographic library. It is also used in the application to validate the license.

The steps to generate a valid and asymmetrically encrypt license using OpenSSL:
\begin{enumerate}
\item Open the private RSA key using {\tt fopen()} and load its contents into memory using {\tt PEM\_read\_RSAPrivateKey()}.

\item Cast the decrypted message from {\tt char*} to {\tt const unsigned char*} as this is what OpenSSL expects as the message's type when encrypting.

\item Encrypt the message using {\tt RSA\_private\_encrypt()}. The method expects the length of the cleartext message, the cleartext message itself, where to store the encrypted message, the private RSA key and the padding of the cleartext message. If the encryption process was unsuccessful, the {\tt bufferSize} will get -1 as its value, otherwise it will get the encrypted messages size.

\item Create a {\tt vector} of {\tt uint8\_t} values out of the encrypted message. This is done because one of the design goals of the backend was to not use pointers between method calls to reduce the risk of memory leaks.

\item Clean up memory and close the open file descriptor to the private key.
\end{enumerate}

\begin{lstlisting}
vector<uint8_t> generateEncryptedMessage(string message, string privateKey) {
    // Step 1
    const char *privkeyPath = (const char*) privateKey.c_str();
    unsigned char *encryptedMessage = new unsigned char[1024]; //should be enough
    FILE *privateKeyFile = fopen(privkeyPath, "r");
    RSA *rsaPrivate = PEM_read_RSAPrivateKey(privateKeyFile, NULL, NULL, NULL);
    if(rsaPrivate == NULL) {
	    RSA_free(rsaPrivate);
	    delete [] encryptedMessage;
    	    fclose(privateKeyFile);
        return vector<uint8_t>();
    }

    // Step 2
    const unsigned char *decryptedMessage = (const unsigned char*) message.c_str();
    
    // Step 3
    int bufferSize = RSA_private_encrypt(message.length(), decryptedMessage,
                                        encryptedMessage, rsaPrivate, RSA_PKCS1_PADDING);
    if(bufferSize == -1) {
        RSA_free(rsaPrivate);
	    delete [] encryptedMessage;
    	    fclose(privateKeyFile);
        return vector<uint8_t>();
    }

    // Step 4
    vector<uint8_t> retVec;
    for(int i = 0; i < bufferSize; i++) {
        retVec.push_back((uint8_t)encryptedMessage[i]);
    }


    // Step 5
    RSA_free(rsaPrivate);
    delete [] encryptedMessage;
    fclose(privateKeyFile);

    return retVec;
}
\end{lstlisting}

To store the encrypted message in a database or to transmit the message through a RESTful web API to the application it is encoded to Base64. The Qt framework provides a simple member method to do that: {\tt QString::toBase64()}. The sample method to decrypt an encrypted message from memory looks similar. The difference is that, on the application side, the public key should not be stored on the persistent storage device of the user as this would allow modifications easily of the public RSA key. Instead, the public RSA key is compiled into the binary of Beat The Score and read from memory. To do this a {\tt static const unsigned char public\_pem[]} is generated out of the public key file before compilation of Beat The Score. This is then used in the decryption method on the client side. {\tt BIO\_new\_mem\_buf} and {\tt PEM\_read\_bio\_PUBKEY} are used in sequence to make the OpenSSL library use the key that is stored inside {\tt public\_pem}. According to the official OpenSSL documentation (https://www.openssl.org/docs/crypto/bio.html), {\tt BIO} in OpenSSL is "an I/O abstraction" and it can be used to "transparently handle SSL connections, unencrypted network connections and file I/O".

\begin{lstlisting}
QString decryptMessage(unsigned char *message, int size)
{
    int pem_size = sizeof(public_pem)/sizeof(*public_pem);
    BIO *publicKeyBuffer = BIO_new_mem_buf((void*)public_pem, pem_size);
    RSA *rsaPublic;
    EVP_PKEY *pkey = PEM_read_bio_PUBKEY(publicKeyBuffer, NULL, NULL, NULL);
    if (pkey) {
        rsaPublic = EVP_PKEY_get1_RSA(pkey);
        if (rsaPublic) {
            RSA_up_ref(rsaPublic);
        }
        EVP_PKEY_free(pkey);
    }
    if(rsaPublic == NULL) {
        return NULL;
    }
    unsigned char *decryptedMessage = new unsigned char[1024];
    int bufferSize = RSA_public_decrypt(size,
                                        message,
                                        decryptedMessage,
                                        rsaPublic,
                                        RSA_PKCS1_PADDING);
    if(bufferSize == -1) {
        return NULL;
    }

    RSA_free(rsaPublic);
    BIO_free(publicKeyBuffer);
    QString retStr((const char*)decryptedMessage);
    delete [] decryptedMessage;
    if(retStr.isEmpty() || retStr.isNull()) {
        return NULL;
    }
    return retStr;
}

\end{lstlisting}